\begin{Arabic}
\chapter*{ملخص}
\addcontentsline{toc}{chapter}{{\arabicfont ملخص}}

\vspace{0.5cm}

تُعدّ عملية ترحيل البيانات من أبرز التحديات التقنية عند نشر نظام تخطيط موارد المؤسسات. في سياق منصة Axelor Open Suite، تعتمد هذه العملية حالياً على برامج Java مخصصة لكل مشروع عميل — وهو نهج مكلف وغير قابل لإعادة الاستخدام وصعب الصيانة.

يقدّم هذا التقرير تصميم وتطوير إطار عمل ETL عام قائم على بوابات القرار، مبني على Apache Hop، يهدف إلى توحيد وأتمتة استيراد البيانات نحو Axelor. يُشكّل هذا الإطار قواعد العمل الخاصة بنظام ERP في تحويلات تصريحية قابلة لإعادة الاستخدام، تغطي أربعة مجالات وظيفية متزايدة التعقيد: المنتجات، الشركاء، أوامر البيع والفواتير.

يقدّم الحل مساهمتين متكاملتين: (1) بنية أنابيب قائمة على بوابات القرار مع حل تلقائي للتبعيات بين الكيانات، تعمل باستقلالية تامة عن واجهة برمجة تطبيقات Axelor؛ (2) نمط ETL بمساعدة الذكاء الاصطناعي يتألف من 3 مهارات ذكاء اصطناعي متخصصة وتطبيق ويب (Axelor Mapper) يؤتمت دورة الاستيراد الكاملة — من التحليل التفاعلي لمخطط المصدر إلى توليد الأنابيب المتوافقة.

\textbf{الكلمات المفتاحية:} ETL، Apache Hop، ERP، Axelor، أنابيب بوابات القرار، ترحيل البيانات، نموذج لغوي كبير، وكيل ذكاء اصطناعي، توليد الشيفرة.

\end{Arabic}
